\section{The stakes of the choice: delta}
\label{sec:ssrdelta}

We have a one-dial family of transports, and a statement from
\cref{sec:ssrmissing} to establish: the dial is a hedge ratio.  This
section derives the delta correction each regime implies and evaluates
it on the frozen node.  Two small reminders before the computation.
\emph{Vega} is an option price's sensitivity to its implied
volatility.  And on this board variance time equals calendar time ---
Chapter~8 ran its clock solver on SPY and found nothing to install ---
so $\tau=t$ throughout, and we write $\tau$.

Start with the object the transport controls directly: the volatility
of a \emph{fixed strike}.  Compose the two rules of
\cref{def:ssrsigns}.  A strike with today's label $k_{0}$ has
tomorrow's label $k_{0}-H$ (quote rule), and the transported curve
evaluated there (curve rule, i.e.\ \eqref{eq:ssrshift}) gives
\begin{equation}
\sigma_{K}^{\mathrm{new}}
=\sigma_{\mathrm{new}}(k_{0}-H)
=\sigma_{\mathrm{old}}\big((k_{0}-H)+H\big)+(\mathcal{R}-1)s_0H
=\sigma_{K}^{\mathrm{old}}+(\mathcal{R}-1)\,s_0\,H .
\label{eq:ssrstrike}
\end{equation}
The re-index and the label slide cancel \emph{exactly} --- no
first-order expansion was needed --- and what is left is the level.  A
fixed strike's volatility responds at the rate $(\mathcal{R}-1)s_0$
per unit log-forward move, at every strike: zero under sticky-strike
(by construction), $-s_0$ under sticky-moneyness, $+s_0$ under
sticky-local-vol.  This is the missing derivative of the chapter's
title, made explicit: choose $\mathcal{R}$ and you have chosen
$\dd\sigma_K/\dd\log F$.

Now push it into the delta.  A dollar call is
$C=D\,F\,\Black(k,w)$, with $\Black$ the normalized Black call of
Chapter~2 (eq.~\ref{eq:black}) and $w=\sigma_K^2\tau$ the strike's
total implied variance; the discount $D$ rides along untouched and we
set it aside, quoting deltas in the forward (Black) convention.
Differentiate along the move.  Three things change at once: the
prefactor $F$, the label $k=\log(K/F)$ (down by $\dd\log F$), and ---
this is the new part --- the volatility, by \eqref{eq:ssrstrike}.  Per
unit of $\dd\log F$, and dividing by $F$ to convert to a delta (since
$\dd F=F\,\dd\log F$),
\[
\Delta_{\mathrm{tot}}
=\Black-\partial_k \Black
+\partial_w \Black\;\frac{\dd w}{\dd\log F}.
\]
Both Black partial derivatives were computed in Chapter~2's appendix
and recalled in Chapter~6, which also performed the collapse the
first two terms need: $\Black-\partial_k\Black=\PhiN(d_+)$, the Black
forward delta.  For the third term, differentiate
$w=\sigma_K^{2}\tau$ and insert \eqref{eq:ssrstrike} together with
the vega formula $\partial_w \Black=\phin(d_+)/(2\sqrt{w})$:
\[
\frac{\dd w}{\dd\log F}
=2\sigma_K\tau\,(\mathcal{R}-1)s_0,
\qquad
\partial_w \Black\;\frac{\dd w}{\dd\log F}
=\frac{\phin(d_+)}{2\sqrt{w}}\;2\sigma_K\tau\,(\mathcal{R}-1)s_0
=\phin(d_+)\sqrt{\tau}\,(\mathcal{R}-1)s_0,
\]
where the last step used $\sqrt{w}=\sigma_K\sqrt{\tau}$: the
volatility cancels.  Together,
\begin{equation}
\Delta_{\mathrm{tot}}
=\PhiN(d_+)+\phin(d_+)\,\sqrt{\tau}\,(\mathcal{R}-1)\,s_0 .
\label{eq:ssrdelta}
\end{equation}
Read it in words: the true delta is the Black delta plus a smile
correction --- the normal density at $d_+$, times the square root of
maturity, times the fixed-strike volatility response the regime
implies.  Between the extreme named regimes the correction swings from
$-\phin(d_+)\sqrt{\tau}\,s_0$ to $+\phin(d_+)\sqrt{\tau}\,s_0$, so the
$\mathcal{R}=0$ versus $\mathcal{R}=2$ gap is
\begin{equation}
\text{delta gap}
=2\,\phin(d_+)\,\sqrt{\tau}\,\lvert s_0\rvert ,
\label{eq:ssrgap}
\end{equation}
a pure product of three quantities the desk already knows.

\begin{figure}[htbp]
\centering
\paperfig{\textwidth}{fig_ssr_delta.pdf}
\caption{One book, three deltas (frozen hero node,
eq.~\eqref{eq:ssrdelta}).  (a)~Total forward call delta across strikes
under the three regimes; at the marked $k=-0.10$ strike --- a
\MacSsrDeltaMarkPutDelta-delta put --- the $\mathcal{R}=0$ and
$\mathcal{R}=2$ answers are \MacSsrDeltaGapPts\ delta points apart.
(b)~The gap \eqref{eq:ssrgap} across strikes: a scaled vega profile,
peaking at \MacSsrDeltaMaxPts\ delta points at the money.}
\label{fig:ssrdelta}
\end{figure}

\Cref{fig:ssrdelta} evaluates both formulas on the frozen node, whose
skew is $s_0=\MacSsrHeroSkew$.  In panel~(a), each curve is
$\Delta_{\mathrm{tot}}$ across strikes for one regime, in delta points
(hundredths).  One ordering holds across the whole axis:
the sticky-moneyness delta is the largest and the sticky-local-vol
delta the smallest --- at every strike, since the correction
$\phin(d_+)\sqrt{\tau}(\mathcal{R}-1)s_0$ carries one sign across the
axis.  The ordering is worth one plain sentence.  Under sticky-local-vol a falling market raises this put's
volatility (eq.~\eqref{eq:ssrstrike} with $\mathcal{R}=2$,
$s_0H>0$ when $H<0$); a long call at that strike therefore loses
less on the way down than Black alone predicts, which is the same as
saying its effective delta is smaller.  At the marked strike --- ten
percent below the money, a \MacSsrDeltaMarkPutDelta-delta put in
Black's own units --- the vertical bracket measures the disagreement:
\MacSsrDeltaGapPts\ delta points between the two outer regimes, on the
same book, against the same fitted surface.  A desk running a hundred
such contracts hedges a difference of \MacSsrDeltaGapPts\ underlying
units, purely on the choice of dial.

Panel~(b) plots the gap \eqref{eq:ssrgap} across strikes, and its
shape should look familiar: $\phin(d_+)$ is (up to normalization) the
vega profile, so the gap is vega rescaled by the skew.  It peaks at
the money --- \MacSsrDeltaMaxPts\ delta points on this node --- and
fades in the far wings, where prices respond too weakly to volatility
for the regime to matter.  Notice what the peak location means: the
regime choice is \emph{most} consequential exactly where the book is
most liquid and the desk feels most confident.  The dashed vertical
marks the panel~(a) strike: at \MacSsrDeltaGapPts\ points it is
nowhere near the worst case.

We now know what the dial costs: double-digit delta points through the
liquid belly of an ordinary index node.  \Cref{rem:ssrident} says the
surface cannot set the dial.  The natural next question is whether
anything else in this book can.  One object can try: the one model of
Part~I that carries every expiry at once, Chapter~4's local-volatility
field.  Holding \emph{it} fixed while spot moves produces a transport
with no dial at all --- that is the next section.
